\documentclass[conference]{IEEEtran}
\IEEEoverridecommandlockouts
\usepackage{cite}
\usepackage{amsmath,amssymb,amsfonts}
\usepackage{algorithmic}
\usepackage{graphicx}
\usepackage{textcomp}
\usepackage{xcolor}
\def\BibTeX{{\rm B\kern-.05em{\sc i\kern-.025em b}\kern-.08em
    T\kern-.1667em\lower.7ex\hbox{E}\kern-.125emX}}
\begin{document}
\title{From Text to Traffic: Repurposing Frozen Generative Models for DDoS Detection\\

}

\author{\IEEEauthorblockN{Naren Medarametla}
\IEEEauthorblockA{\textit{School of Computer Science Engineering} \\
\textit{Vellore Institute of Technology}\\
Chennai, India \\
naren.medarametla2023@vitstudent.ac.in}
\and
\IEEEauthorblockN{Rajnandini}
\IEEEauthorblockA{\textit{School of Computer Science Engineering} \\
\textit{Vellore Institute of Technology}\\
Chennai, India \\
email address or ORCID}
}

\maketitle

\begin{abstract}
Distributed Denial-of-Service (DDoS) attacks remain a critical cybersecurity threat, aiming to disrupt services by overwhelming 
target networks with volumetric traffic floods. While machine learning models are increasingly deployed for Network Intrusion Detection 
Systems (NIDS), they traditionally struggle with zero-shot generalization, frequently failing to detect novel, out-of-distribution 
attack protocols due to rigid global feature scaling. Taking inspiration from the DoLLM framework, we propose an adapted architecture 
that repurposes a frozen Large Language Model for cross protocol network traffic analysis. Our approach pairs a 1D-Convolutional temporal 
flow tokenizer and Multihead Attention pooling with a localized sequence-level normalization technique. By normalizing network flows 
strictly relative to their immediate temporal windows rather than a global dataset baseline, the model achieves scale-invariant reasoning 
to successfully identify the structural physics of malicious bursts. Evaluated on completely unseen attack protocols, our adapted 
architecture demonstrates robust cross protocol generalization, achieving a zero-shot F1 score of 0.9922.
\end{abstract}

\begin{IEEEkeywords}
Deep Learning, Large Language Models, Distributed Denial-of-Service, Network Security, Network Intrusion Detection Systems
\end{IEEEkeywords}

\section{Introduction}

\section{Literature Review}

\section{Methodology}

\section{Results}

\section{Discussions}

\section{References}

% \begin{thebibliography}{00}

% \end{thebibliography}
\vspace{12pt}

\end{document}
